\chapter{多因子模型、估计与研究答案}

\begin{enumerate}[leftmargin=*]
  \item 预测型横截面回归写成 $r_{i,t+1}=a_t+b_t x_{i,t}+u_{i,t+1}$，回答的是给定样本和时点的预测关系；经典风险价格来自 $r_{i,t}=\beta_i^T\lambda+\alpha_i+\varepsilon_{i,t}$ 的定价结构；因果效应还需处理、潜在结果和识别假设。显著的 $b_t$ 不能单独证明 $\lambda$ 或因果。
  \item Pearson 使用原始距离和协方差，极端值会改变结果；Rank IC 把值换成秩，只看次序。对单调但弯曲关系，Rank IC 可能高而 Pearson 较低；对少数离群点，Pearson 会翻转而 Rank IC 较稳。并列值的 tie 规则必须固定。
  \item 两期斜率均值为 $\bar b=(0.020+0.030)/2=0.025$。样本离均差为 $-0.005,+0.005$，样本方差为 $[(.005)^2+(.005)^2]/(2-1)=5\times10^{-5}$，均值标准误为 $\sqrt{5\times10^{-5}/2}=0.005$。若有 $T$ 期，应按时间维度而非资产维度估计误差。
  \item 第一遍对每个资产 $i$ 用时间序列回归 $r_{i,t}=\alpha_i+\beta_i^Tf_t+\varepsilon_{i,t}$，保存 $\widehat\beta_i$；第二遍在每个 $t$ 回归 $r_{i,t}$ 对 $\widehat\beta_i$，得到 $\widehat\lambda_t$，再取时间均值。因为第二遍回归量是估计的 beta，存在生成回归量误差和估计误差相关，不能直接套简单 OLS 标准误。
  \item 若所有资产都有相同 alpha，第二遍截距吸收它，风险价格保持不变；若 alpha 与 beta 相关，横截面斜率和 beta 排序都被污染。把 alpha 改成指定值后，先重算每个资产平均收益，再重新跑第二遍并比较截距、斜率和标准误，不能只改最终表格。
  \item 相关候选使 BH 的独立/正相依条件需要额外审查。保存完整候选族的相关矩阵和所有 p 值，按预先协议选择 BH、BY、置换或 block 重采样；把筛出后 20 个再做 BH 会改变检验族，不能声称控制原始 FDR。
  \item 预测系数摘要缺少可交易定义、beta 稳定性、定价模型、生成回归量误差、成本后样本外和多重检验协议。逐项把“对象—假设—证据—失败条件”写入审计表，不能用一个显著性星号替代。
  \item 最终期一旦参与信号尺度、窗口、特征或检验族选择，就条件化在选择事件上。把选择移到内层并冻结全部规则，再在从未查看的后移测试期运行一次；若没有新测试期，只能把结果标为探索性。
\end{enumerate}

\subsection*{基础衔接：独立计算}
均值 $0.025$；平方偏差和为 $2(0.005)^2$，标准误为 $\sqrt{2(0.005)^2/(2\cdot1)}=0.005$。只有两期，不能据此期待渐近正态近似可靠。
